\documentclass[final]{clv2025}

\jvol{vv}
\jnum{nn}
\jyear{2025}
%There is no need for the authors to change the above.

\dochead{An Example of Document Head} 
% Possible options: Long Paper, Short Paper, Survey, Survey Proposal, Position Paper, Last Words, Book Review, Dissertation Award, Lifetime Achievement Award, Squibs and Discussions, Featured Article

\pageonefooter{Action editor: \{action editor name\}. Submission received: DD Month YYYY; revised version received: DD Month YYYY; accepted for publication: DD Month YYYY.}
% To be filled when the paper is accepted. The editorial office will share the relevant information with the authors. Example: Action Editor: Wei Lu. Submission received: 10 February 2023; revised version received: 30 May 2023; accepted for publication: 15 June 2023.

\usepackage{amsmath}
\usepackage{booktabs}

\runningtitle{Running Article Title Here}
\runningauthor{Author's Surnames Here}

\begin{document}

\title{This is an Example of Article Title}

\author{Author A\thanks{Equal contribution}$^{,1}$, Author B$^{*,1}$, Author C$^{*,1}$, Author D$^{1}$, Author E$^{2}$, Author F$^{1}$, Author G$^{1}$, Author H$^{2}$, Author I$^{1}$, Author J$^{3}$, Author K\thanks{Corresponding authors}$^{,2}$, Author L$^{**,1}$}

\affilblock{
    \affil{Affiliation 1\\\quad \email{corresponding@author\_L.email}}
    \affil{Affiliation 2\\\quad \email{corresponding@author\_K.email}}
    \affil{Affiliation 3}
}

\maketitle

\begin{abstract}
This is example text for article abstract. This shouldn't be treated as a live paper. This document provides a very brief overview of the file as well as some general instructions on coding or tagging of different elements of a document. This is just a template for author to follow the uniform tex tagging for the environments. This document just gives you some brief instructions on using this package file. Most common formatting commands are included in this document.
\end{abstract}

\section{This is an Example of A Head}

The text of this document is taken from some essays and stories available in public domain on the internet, here we are not concerned with the data of this document while this should be treated as sample on dummy text.

\subsection{This is an Example of B Head}

I love to read. I love books. I am surrounded by books, thousands of them. The sight of them pleases and soothes me. Pulitzer Prize-winners in first editions, two entire bookcases of movie books, a shelf on religion with every major holy book (including B'hai), five 7-foot cases of fiction and non-fiction alphabetical by author, humor, drama, World Book Encyclopedias, an unabridged dictionary, travel books, almanacs, reference books, many books of quotes, atlases, the complete hardcover Yale Shakespeare.

\subsubsection{This is an Example of C Head}

The first book I remember giving me a feeling C Head of cozy, warm, completely transported pleasure was ``The Godfather'' in 1969 when I was eleven years old. I recall fondly curling up on my bed and reading in complete absorption night after night. That such things occurred in the world, and that I should be allowed to know about them was just way too wonderful for words.

\paragraph{This is an Example of D Head}

I've always tried to read them between other kinds of books because they seem like such brain candy to me. Interestingly, I think, I don't read much fiction anymore. I used to only read fiction and science fiction when I was younger. My tastes have changed. Now, the most luxurious kind of book for me is a long biography or autobiography. I recently read ``Wizard: The Life \& Times of Nikola Tesla'' and completely enjoyed every moment.
\begin{enumerate}
\item And then there's all of the filmmaking and film history books I've been devouring endlessly since I was a child.
\item My favorite biographer is Robert K. Massie and his books ``Nicholas \& Numbered Alexandra,'' ``Peter the Great'' and ``Dreadnought'' are as good as books get List from my limited perspective.
\item That it takes Mr. Massie ten years a book is the difficult part for his fans. This is the component of a theory that deals with considerations about the quality of descriptions.
\end{enumerate}
Many people's first question upon seeing all of my books is, ``Have you read them all?'' No, of course not. I don't think I've read half of them (most of the books I've read I took out of the library). Part of the great pleasure is being surrounded by so many unknown adventures.

\begin{itemize}
\item The information it acquires will be described at a semantically {\em appropriate} level of generalisation.
\item It will be defined as a knowledge-poor and unsupervised strategy.
\end{itemize}
It's still absolutely incredible and wonderful to me to be transported back to London in the 1890s as little Charlie and his older brother Sydney come home to find that their mother has gone crazy and has been taken away to an insane asylum. They both end up in a workhouse, just like Oliver Twist.\footnote{Chaplin wrote the book in 1964 (when I was six), which is not all that long ago for someone to have grown up living in a Charles Dickens novel.}

Good books and movies and stories in general transport us to different times and places and let us see how other people have dealt with this dilemma called life, which, let's face it.
\begin{displaytext}
Display(test$_{000}$,example$_{xyzx}$,text$_{000}$),
\end{displaytext}
Right now I'm reading Charlie Chaplin's autobiography, a book I've owned for 20 years and which I am finally reading. Right now I'm reading Charlie Chaplin's autobiography, a book I've owned for 20 years and which I am finally reading.
\begin{displaytext}
[test] {\em for} [display] {\em text},
\end{displaytext}
we none of us will ever get a particularly good enough handle on. But it's always good to see how others dealt or tried to deal with it. We can even occasionally pick up useful ideas.

Books are also a way to set standards and goals for ourselves. Do you think you're really smart and an over-achiever? Read about Tesla or Edison or Alexander the Great. You think you're a filmmaker? Read what Orson Welles or Alfred Hitchcock had to say. Do you think you're a good writer? Read Willa Cather or Joseph Conrad or Norman Mailer and see what's actually possible.
\begin{extract}
Here's a few to start with. . .When I began my college career, in August of 2000, I was amazed at the number of freshmen who, like me, wanted to major in Pre-Medicine. More than half of the freshman class, it seemed, would become doctors some day. However, as time marched on, the numbers began to decline. As different interests, less-than-stellar grades, and various other influences came into our lives, it became evident that a smaller number of students would actually still be pre-med when it came time to take the MCAT. [p. 100]
\end{extract}
Everything we humans know we've gone to the trouble of writing down in books, and all you have to do is pick them up and read them.\footnote{Everything we humans know we've gone to the trouble of writing down in books, and all you have to do is pick them up and read them:
\begin{align}
N(x|y)=\sum_{i=1\dots n}N(x_i|\Gamma(y_1\cdots y_{i-1},x)) \nonumber
\end{align}
This is just nice.
}

The rigors of the MCAT narrowed the numbers a bit more, and now, as we satisfy our last requirements and prepare to graduate, only about 15\% of my graduating class will continue with the intent of becoming physicians:
\begin{align}
\Gamma=\left\{
\begin{array}{ll}
x&\quad \text{if $x$ is variable $\left<X\rightarrow Z\right>$}\\
y&\quad \text{otherwise}
\end{array}
\right.
\nonumber
\end{align}
I am among them. Even though my college experience has allowed me to research many different possible life paths, my desire to be a physician has only been affirmed and strengthened these past 4 years, where
\begin{align}
P(b,\beta)
=\sum_{l=0}^n \alpha_a \kappa(x) = \beta\cdot \Phi(x)
\end{align}

In college I chose to join my university's see Equation (1). Co-ed service fraternity, Alpha Phi Omega. Through serving the local community in several different capacities with my brothers and sisters, I have learned that I am called to a service vocation. I love to help others in whatever way I am able, and this knowledge affirms my desire to help my fellow man as a healer.
\begin{align}
\log N(x|s_i)&=\sum_i \log \frac{e^{F(b,a)}}{\alpha}\nonumber
\\
j&=1^{n_i}e^{F(i_j,b_a)}\nonumber
\\
k&=\frac{M(x_i)}{\alpha_0}m(i_j)\nonumber
\end{align}
I am also a musician at heart. I was involved with the Baylor School of Music each year of my undergraduate career as a member of the Golden Wave Marching Band and the Concert Band.
\begin{align}
P_i(\eta)=x_i N(x_i,\alpha)
\end{align}
Music is certainly a love of mine, and I definitely enjoyed being able to express myself
in this manner; but I also discovered that I love medicine more. 
\begin{algorithm}
\item[Step 1] Through this experience I was encouraged that I made the right decision in pursuing medicine instead of music.
\item[Step 2] Since June 2002, I have worked as an Emergency Medical Technician at Scott and White Hospital in Temple, Texas.
\end{algorithm}
It has most definitely done so. Every time I go to work, I am thankful for the opportunities it affords me. I have been able to experience one perspective of medicine (that is, medicine outside of the hospital) that many doctors do not often see. I have already had many patients, and have become very comfortable with patient care. Whereas many of my fellow medical students will have to learn how to take vital signs and conduct a cursory assessment of a patient, I will have been doing so for several years. My job has not discouraged me from pursuing medicine in any way; it has only affirmed my decision and given me a taste of something for which I am passionate. (see Figure \ref{fig:ex}) and other experiences that have composed my undergraduate career have absolutely confirmed that my passion, my heart, is for medicine. I honestly believe that my calling is to serve others as a healer, not only of the body, but of the mind, soul, and spirit.

\begin{figure}[t!]
\begin{center}
\includegraphics[width=0.25\textwidth]{ex.pdf}
\end{center}
\caption{This is an example of figure caption.}
\label{fig:ex}
\end{figure}

\begin{table}[b!]
\caption{This is an example of table caption}
\label{tab:ex}
\begin{tabular}{llll}
\toprule
Head&\multicolumn{3}{c}{Spanner Head}\\
\cmidrule{2-4}
&Head2&Head3&Head4\\
\midrule
Entry One&07.5&83.4&100\\
Entry Two&80.5&80.7&92\\
Entry Three&90.3&90.4&74\\
\bottomrule
\end{tabular}
\end{table}

Table \ref{tab:ex} the doctor's office. Most kids put a heating pad on their forehead so they could fake illness and stay home from school. When I'd have a fever, I'd use a damp cloth to cool myself down so I could go. I knew it shouldn't be that way. I shouldn't be afraid of the doctor. My goal is to make the experience for children enjoyable. My office will be different. Beanbag chairs instead of couches. The Disney Channel instead of CNN. Coloring books instead of Newsweek. All patients will be given a prescription for candy when they walk in with the instructions to ``Take two M\&M's and call me in the morning.'' The focus will be on the child from the moment he or she walks in the office.

\appendix

\appendixsection{This is an Example of Appendix Section}

I fear that too often, people are given their certification solely because they can pass a test and make it through the necessary training. I know many premed students who only speak science, which may be the language among doctors, but often intimidates and confuses the patient.
\begin{align}
\Gamma=\left\{
\begin{array}{ll}
x&\quad \text{if $x$ is variable $\left<X\rightarrow Z\right>$}\\
y&\quad \text{otherwise}
\end{array}
\right.
\nonumber
\end{align}
While I do feel that a love and respect for science is paramount in the development of a good doctor, I feel that all too often the love and respect of the patients is lost within this quest.
\begin{align}
P(b,\beta)
=\sum_{l=0}^n \alpha_a \kappa(x) = \beta\cdot \Phi(x)
\end{align}
While having scored in the 97th percentile on the MCAT coupled with my performance in all of the math and science classes I have taken shows that I have the necessary for understanding of science required for a career in medicine.
\begin{align}
P_i(\eta)=x_i N(x_i,\alpha)
\end{align}
It is my passion for communication and people that will shine through me as a doctor.

\begin{figure}[h!]
\begin{center}
\includegraphics[width=0.25\textwidth]{ex.pdf}
\end{center}
\caption{This is another example of figure caption.}
\label{fig:ex2}
\end{figure}

Figure \ref{fig:ex2} is here. That day. . .that smile, has taught me that a career in pediatric medicine is not merely a career in science, it is a career in people. It is imperative to remember that while in the lab I am dealing with the illness of a child, in the room I am dealing with a child who is ill. It is my hope that this belief is not unique to me, but is universal to all -- for it is based on this belief that a doctor's office evolves from a place that is dreaded by patients, to one that is embraced by children.

\appendixsection{This is an Example of Appendix Section}

\begin{table}[h!]
\caption{This is another example of table caption}
\label{tab:ex2}
\begin{tabular}{llll}
\toprule
Head&\multicolumn{3}{c}{Spanner Head}\\
\cmidrule{2-4}
&Head2&Head3&Head4\\
\midrule
Entry One&07.5&83.4&100\\
Entry Two&80.5&80.7&92\\
Entry Three&90.3&90.4&74\\
\bottomrule
\end{tabular}
\end{table}

This same balance is something that I have striven to attain with my involvement in the community. As you can see from Table \ref{tab:ex2}, ather then focusing exclusively on volunteering that is centered on medicine, I have tried to spend as much time as possible working with the medium of the pediatric field - kids.
\begin{align}
P(b,\beta)
=\sum_{l=0}^n \alpha_a \kappa(x) = \beta\cdot \Phi(x)\label{eq:a}
\end{align}
Since the age of eight I have done magic professionally \cite{article1} and have been able to use this talent to entertain children \cite{article2,article3}. This is clear from Equation \ref{eq:a}. The true magic is not in correctly identifying the card that was chosen, but in being able to momentarily suspend belief for a child; being able to make them feel that anything is possible. After I can convince them of this, it is much easier for them to accept that they too can defy the impossible. This interaction with children, which has led to a heightened understanding of how they think, is something that cannot be taught in a classroom. As \citet{,article4,article5,article6} mentioned, it is an education that does not stop with the framing of a diploma. It continues daily.

\begin{acknowledgments}
This is an example of acknowlwdgment. This is a sample paper presented jost for the coding of different elements of a document in \LaTeX\ using this package file. We thank all readers.
\end{acknowledgments}

%\newpage
%If you would like to use \starttwocolumn below for references: Unfortunately, there is a known bug in the existing \starttwocolumn command when used without \newpage above it. If you find issues such as whole pages disappearing, add \newpage above \starttwocolumn as a temporary fix. We will provide a permanent solution in a future update.
%\starttwocolumn

\bibliographystyle{compling}
\bibliography{COLI_template}

\end{document}
